% docs/physics/latex/08_sidis.tex - SIMC++ physics note 08: semi-inclusive (e,e'h)X production (SIDIS).
% Standalone PDF docs/physics/pdf/SIMCpp_Physics_08_sidis.pdf (copied to src/reactions/sidis/ and src/physics/sidis/);
% a chapter of the book SIMC++ Physics Reference. Rules: docs/physics/README.md; structure: template_reaction.tex.
% Every display cites the C++ function first and the Fortran (upstream 59ccdf6) line second.
\documentclass{simcnote}
\simccommon
\begin{document}
\simcnotemeta{08}{SIDIS (e,e\pr h)X}{Semi-inclusive production (e,e\pr\texorpdfstring{$\pi^\pm$}{pi+-})X, (e,e\pr\texorpdfstring{$\pi^0$}{pi0})X and (e,e\pr K\texorpdfstring{$^\pm$}{+-})X}%
  {Reaction plug-in \texttt{sidis\_production} with the physics model \texttt{sidis}: CTEQ5M parton distributions,
   Bosted pion and fDSS kaon fragmentation, Christy inclusive structure functions}
\simcnotemaketitle
\simcabstract{Semi-inclusive electroproduction of a pion or a kaon on hydrogen, deuterium or a heavier nucleus, with the
electron and the hadron detected in coincidence and the rest of the final state unobserved. Both final energies and
both directions are generated, nothing is solved from the kinematics. The weight is a factorised leading-order model:
the inclusive cross section (from the Christy--Bosted structure-function fit) times the hadron multiplicity
$\sum_q e_q^2q(x)D_q^h(z)/\sum_q e_q^2q(x)$ (CTEQ5M parton distributions; P.~Bosted's fit of the pion fragmentation
functions, fDSS for kaons), times a Gaussian transverse-momentum distribution and a flat azimuth. Notation, glossary and
common physics: note 00.}
\simcnotetoc

% =====================================================================================================================
\section{The reaction in SIMC++}\lbl{sec:reaction}

\subsection{Process and scope}
The reaction is $e + A \to e' + h + X$ with $h=\pi^+,\pi^-,\pi^0,K^+,K^-$, on hydrogen, deuterium or a nucleus with
$A\ge3$. What is simulated: the full event chain of note 00 (beam, energy loss, multiple scattering, radiation on the three
legs, both spectrometers, decay of the hadron in flight or its survival probability, the two photons of a $\pi^0$), with a
factorised cross section built from parton distributions and fragmentation functions. What is not simulated: any
azimuthal modulation (the $\varphi_{pq}$ dependence is flat, no Cahn, Boer--Mulders or higher-twist terms), target-mass
corrections other than the scaling variable of the pion fragmentation fit, exclusive and resonance channels (they are
cut by two missing-mass thresholds, \secr{sec:complete} and \secr{sec:xsec}), the unobserved hadrons, Fermi motion in
nuclei with $A\ge3$ (only deuterium has it, and only in the cross section), polarised targets in the cross section, and
any nuclear effect beyond the EMC ratio of the inclusive part and an attenuation factor of the multiplicity.

\subsection{Components}
\begin{simctable}{footnotesize}{L{0.14\linewidth}L{0.16\linewidth}L{0.15\linewidth}L{0.25\linewidth}L{0.30\linewidth}}
\caption{The SIMC++ components of SIDIS (one component per row).}\lbl{tab:components}\\
\toprule\rowcolor{simchead}\simcth{role} & \simcth{interface} & \simcth{registered name} & \simcth{class} & \simcth{file}\\\midrule\endfirsthead
\toprule\rowcolor{simchead}\simcth{role} & \simcth{interface} & \simcth{registered name} & \simcth{class} & \simcth{file}\\\midrule\endhead
\bottomrule\endlastfoot
reaction & \code{IReaction} & \code{sidis\_production} & \code{SidisProduction} & \file{reactions/sidis/SidisProduction.cpp}\\
meson base class & (base class) & -- & \code{PionProduction} & \file{reactions/pion/PionProduction.cpp}\\
event completion and reconstruction & (base class) & -- & \code{EepReactionBase} & \file{reactions/EepReactionBase.cpp}\\
cross section & \code{IPhysicsModel} & \code{sidis} & \code{sidis::SidisModel} & \file{physics/sidis/SidisModel.cpp}\\
cross-section kernel & (function) & -- & \code{sidis::peepiX} & \file{physics/sidis/SidisCrossSection.cpp}\\
EMC ratio and hadron attenuation & (functions) & -- & \code{sidis::slacemcfit}, \code{sidis::rhadGlobal} & \file{physics/sidis/SidisCrossSection.cpp}\\
parton distributions & (class) & -- & \code{sidis::Cteq5Pdf} & \file{physics/sidis/Cteq5Pdf.cpp}\\
kaon fragmentation & (class) & -- & \code{sidis::FdssGrid} & \file{physics/sidis/Fdss.cpp}\\
inclusive structure functions & (functions) & -- & \code{sidis::f1f2in21}, \code{sf}, \code{resmod} & \file{physics/sidis/F1F2In21.cpp}\\
dilution factor (not called) & (function) & -- & \code{sidis::semiDilution} & \file{physics/sidis/SemiDilution.cpp}\\
data tables & (files) & -- & \code{cteq5m.tbl}, \code{KANLO.GRID} & \file{data/sidis/}\\
deuteron momentum distribution & (file) & -- & \code{ReadMomentum\-Distribution} (\code{deut.dat}) & \file{setup/SetupTheory.cpp}\\
$\pi^0\to\gamma\gamma$ & (function) & -- & \code{pion::pizeroDecay} & \file{physics/pion/PizeroDecay.cpp}\\
beam, target, radiation, arms, outputs & as in note 00 & \code{standard}, \code{hallc\_cryo}, \code{simc\_rad}, \ldots & & \ctabr{tab:plugins}\\
\end{simctable}

\subsection{Selection and override}
The pipeline selects \code{sidis\_production} with the model \code{sidis} when \code{doing\_semi} = 1
(\ctabr{tab:defaults}). \code{doing\_semi} then takes priority over the meson flags, which only choose the hadron:
\code{doing\_pion} = 1 gives pions, \code{doing\_kaon} = 1 kaons; \code{doing\_hplus} chooses the charge and
\code{doing\_pizero} = 1 (with \code{doing\_pion}) the neutral pion. The target class follows $\mathrm{nint}(A)$: 1, 2,
$\ge3$. The set-up uses $M_h$ = $M_{\pi^\pm}$, $M_{\pi^0}$ or $M_K$, a struck nucleon $M_N=M_p$ and a ``recoil'' of mass
$M_p$ for every target and hadron (the code comment says ``for now, just assuming proton mass''), so that $M_N$ and the
recoil mass cancel in the energy budget. \code{do\_fermi} is switched off (with a warning) on hydrogen. The model
refuses a set-up without \code{doing\_semi}, without a hadron choice, or with \code{doing\_phsp} (note 09). JSON keys:

\begin{simctable}{footnotesize}{L{0.33\linewidth}L{0.33\linewidth}L{0.34\linewidth}}
\caption{JSON keys of this reaction and its model (one key per row).}\lbl{tab:json}\\
\toprule\rowcolor{simchead}\simcth{key} & \simcth{effect} & \simcth{example}\\\midrule\endfirsthead
\toprule\rowcolor{simchead}\simcth{key} & \simcth{effect} & \simcth{example}\\\midrule\endhead
\bottomrule\endlastfoot
\code{pipeline.reaction} & the reaction & \code{"sidis\_production"}\\
\code{pipeline.physics} & the cross-section model & \code{"sidis"}\\
\code{pipeline.physics.data\_dir} & directory of \code{cteq5m.tbl} and \code{KANLO.GRID} & \code{\{"type": "sidis", "data\_dir": "/path"\}}\\
\code{pipeline.physics.verbose} & print the notes on ignored inputs and the model messages & \code{\{"type": "sidis", "verbose": true\}}\\
\end{simctable}
A replacement model must return $\mathrm{d}^5\sigma/(\mathrm{d}\Omega_e\,\mathrm{d}E'\,\mathrm{d}\Omega_h\,\mathrm{d}E_h)$
in $\mu$b/(MeV$^2$\,sr$^2$) per nucleus (\secr{sec:xsec}).

\subsection{Event flow}
\begin{simctable}{footnotesize}{L{0.05\linewidth}L{0.33\linewidth}L{0.62\linewidth}}
\caption{One try of SIDIS (one step per row; the general chain is \ctabr{tab:chain}).}\lbl{tab:flow}\\
\toprule\rowcolor{simchead}\simcth{\#} & \simcth{C++ call} & \simcth{what is specific to this reaction}\\\midrule\endfirsthead
\toprule\rowcolor{simchead}\simcth{\#} & \simcth{C++ call} & \simcth{what is specific to this reaction}\\\midrule\endhead
\bottomrule\endlastfoot
1 & \code{StandardBeam::generate} & nothing\\
2 & \code{SidisProduction::generate} & both slope pairs, the hadron energy, then the electron energy (\eqsr{eq:genhad}{eq:genel}); deuterium: a Fermi momentum (\eqr{eq:fermi})\\
3 & \code{SidisProduction::completeEv} & no solution: $\vec q$, $W$, $\varepsilon$, $\theta_{pq}$, $t$, $\varphi_{pq}$, $\vec P_m$, $E_\mathrm{miss}$; the missing-mass threshold, $z$, $p_T^2$ (\eqsr{eq:mesonkin}{eq:zpt}); the $\pi^0$ photons; two solid-angle jacobians\\
4 & \code{SimcRadiator::radiate} & tail-1 limit from the summed-energy limit (\eqr{eq:radlimits}); a tail-1 photon calls \code{completeEv} again\\
5 & \code{MonteCarlo::run} & decay of $\pi^\pm$/K$^\pm$ in flight when \code{doing\_decay} = 1; for a $\pi^0$ the photons go to a calorimeter arm\\
6 & \code{EepReactionBase::completeRecon} & $E_m^\mathrm{rec}$, missing masses, $t^\mathrm{rec}$, $z^\mathrm{rec}$, $p_T^{2,\mathrm{rec}}$ (\eqr{eq:recon})\\
7 & \code{SidisProduction::completeMain} & $S=1$, $n_\mathrm{tgt}=1$; \code{SidisModel::sigccMain} (\secr{sec:xsec}); Coulomb factor; survival probability for kaons without decay\\
\end{simctable}

% =====================================================================================================================
\section{Kinematics, generated variables and generation volume}\lbl{sec:gen}
\subsection{Generated variables}
\begin{simctable}{small}{L{0.35\linewidth}L{0.65\linewidth}}
\caption{Generated variables (one variable per row; $\xi_i$ uniform random numbers, in this order).}\lbl{tab:generated}\\
\toprule\rowcolor{simchead}\simcth{variable} & \simcth{how}\\\midrule\endfirsthead
\bottomrule\endlastfoot
electron slopes $y'_e$, $x'_e$ & flat in the electron edges (as in note 01)\\
hadron slopes $y'_h$, $x'_h$ & flat in the hadron edges\\
hadron energy $E_h$ & flat, drawn first (\eqr{eq:genhad})\\
electron energy $E'$ & flat in its full generation range (\eqr{eq:genel})\\
Fermi momentum $\vec p_F$ (deuterium only) & from the deuteron momentum distribution, isotropic (\eqr{eq:fermi})\\
\end{simctable}
After the four slopes the hadron energy is drawn in the part of its generation range compatible with the summed-energy
limits $\Sigma$, then the electron energy in its own range:
\begin{feqa}{reactions/sidis/SidisProduction.cpp: SidisProduction::generate}{event.f:289-297}
E_h^\mathrm{lo} &= \max\bigl(E^\mathrm{gen}_{h,\mathrm{min}},\ \Sigma_\mathrm{min}-E'^\mathrm{gen}_\mathrm{max}\bigr),\qquad
E_h^\mathrm{hi} = \min\bigl(E^\mathrm{gen}_{h,\mathrm{max}},\ \Sigma_\mathrm{max}-E'^\mathrm{gen}_\mathrm{min}\bigr),\nonumber\\
E_h &= E_h^\mathrm{lo}+\xi_5\,(E_h^\mathrm{hi}-E_h^\mathrm{lo}),\qquad P_h = \sqrt{E_h^2-M_h^2},\qquad
w_\mathrm{gen}\leftarrow w_\mathrm{gen}\,\frac{E_h^\mathrm{hi}-E_h^\mathrm{lo}}{E^\mathrm{gen}_{h,\mathrm{max}}-E^\mathrm{gen}_{h,\mathrm{min}}},
\lbl{eq:genhad}
\end{feqa}
\begin{feq}{reactions/sidis/SidisProduction.cpp: SidisProduction::generate}{event.f:300-316}
E' = E'^\mathrm{gen}_\mathrm{min}+\xi_6\,(E'^\mathrm{gen}_\mathrm{max}-E'^\mathrm{gen}_\mathrm{min}),\qquad
P'_e=E',
\lbl{eq:genel}
\end{feq}
the try failing when a range is empty. Unlike A(e,e\pr p), the hadron window does not depend on the drawn electron energy
(it is built from the fixed limits of the other energy), so the factor of \eqr{eq:genhad} is the same for every event,
and the electron factor is 1. The slopes are converted to angles with \ceqr{eq:physang}.

For deuterium a Fermi momentum is drawn in every try, whether or not \code{do\_fermi} is set (four random numbers), from
the tabulated momentum distribution $p_k$ with normalised cumulative probabilities $C_k$ (file \code{deut.dat}, $n$
points):
\begin{feqa}{reactions/sidis/SidisProduction.cpp: SidisProduction::generate}{event.f:339-367}
k &= \min\{k: \xi_7\le C_k\}\ (\le n),\qquad p_F = p_\mathrm{lo}+\xi_8\,(p_\mathrm{hi}-p_\mathrm{lo}),\nonumber\\
p_\mathrm{lo} &= \tfrac12(p_{k-1}+p_k)\ (0\text{ for }k=1),\qquad p_\mathrm{hi} = \tfrac12(p_k+p_{k+1})\ (p_n\text{ for }k=n),\nonumber\\
\hat p_F &= (\sin\theta_F\cos\varphi_F,\ \sin\theta_F\sin\varphi_F,\ \cos\theta_F),\quad \cos\theta_F = 2\xi_9-1,\quad
\varphi_F = 2\pi\xi_{10},\nonumber\\
E_m &= M_p+M_n-M_d,\qquad m_s = M_A-M_N+E_m\ (=M_n),\qquad E_F = M_A-\sqrt{m_s^2+p_F^2},
\lbl{eq:fermi}
\end{feqa}
$E_F$ being the off-shell energy of the struck nucleon with the spectator on shell. Without Fermi motion (hydrogen, and
every nucleus with $A\ge3$) $p_F=0$ and $E_F=M_N$. The Fermi momentum does not change the kinematics of the event (the
hadron is generated freely); it enters only the cross section (\secr{sec:fermixs}).

\subsection{Limits}
The vertex missing energy and momentum, the residual mass and kinetic energy are not constrained (all limits 0). With
$M_N$ equal to the recoil mass, energy conservation gives
\begin{feqa}{setup/SetupLimits.cpp: BuildLimits}{init.f:359-363, 411-425, 450-454, 467-469}
E_\gamma^\mathrm{tot,max} &= E_b^\mathrm{max} - E'_\mathrm{min} - E_{h,\mathrm{min}},\nonumber\\
\Sigma_\mathrm{max} &= E_b^\mathrm{max}+d_E,\qquad \Sigma_\mathrm{min} = \max\bigl(0,\ E'_\mathrm{min}+E_{h,\mathrm{min}}-d_E\bigr),
\lbl{eq:sumegen}
\end{feqa}
($E_\gamma^\mathrm{tot,max}=0$ without radiation, \code{Egamma\_gen\_max} when set; each enabled tail gets
$E^\mathrm{max}_{\gamma i}=E^\mathrm{tot,max}_\gamma$), and the generation ranges of the two energies
\begin{feqa}{setup/SetupLimits.cpp: BuildLimits}{init.f:481-488, 504-513}
[E'^\mathrm{gen}_\mathrm{min},E'^\mathrm{gen}_\mathrm{max}] &= [\Sigma_\mathrm{min}-E_{h,\mathrm{max}}-E^\mathrm{max}_{\gamma3},\
\Sigma_\mathrm{max}-E_{h,\mathrm{min}}]\cap[E'_\mathrm{min},\ E'_\mathrm{max}+E^\mathrm{max}_{\gamma2}],\nonumber\\
[E^\mathrm{gen}_{h,\mathrm{min}},E^\mathrm{gen}_{h,\mathrm{max}}] &= [\Sigma_\mathrm{min}-E'_\mathrm{max}-E^\mathrm{max}_{\gamma2},\
\Sigma_\mathrm{max}-E'_\mathrm{min}]\cap[E_{h,\mathrm{min}},\ E_{h,\mathrm{max}}+E^\mathrm{max}_{\gamma3}],
\lbl{eq:genranges}
\end{feqa}
with $E'_\mathrm{min/max}$, $E_{h,\mathrm{min/max}}$ the edges of \ceqr{eq:edges}. With radiation the upper hadron limit
$E_{h,\mathrm{max}}+E^\mathrm{max}_{\gamma3}$ is large, so the hadron is generated far above the hadron acceptance
(relative momenta up to about +100\,\% at typical Hall C settings); most such events fail, which lowers the generation
efficiency of radiative runs.

\subsection{Generation volume}
\begin{feq}{reactions/EepReactionBase.cpp: EepReactionBase::genVolume}{simc.f:383-401}
V_\mathrm{gen} = \Delta\Omega_e\,\Delta\Omega_h\,\bigl(E'^\mathrm{gen}_\mathrm{max}-E'^\mathrm{gen}_\mathrm{min}\bigr)
\bigl(E^\mathrm{gen}_{h,\mathrm{max}}-E^\mathrm{gen}_{h,\mathrm{min}}\bigr)\quad[\mathrm{sr^2\,MeV^2}],
\lbl{eq:genvol}
\end{feq}
the six-fold volume of \ceqr{eq:genvol}.

% =====================================================================================================================
\section{Completing the event}\lbl{sec:complete}
\subsection{Vertex}
Nothing is solved: $\nu$, $Q^2$, $\vec q$ and $x$ follow from \ceqsr{eq:nuq2}{eq:xbj} with the generated $E'$, and the
hadron keeps its generated energy and direction. The meson block (shared with the pion reaction) computes, with the
struck-nucleon mass $M_N$,
\begin{feqa}{reactions/pion/PionProduction.cpp: PionProduction::pionMainKinematics}{event.f:723-732}
W^2 &= M_N^2+2M_N\nu-Q^2,\qquad W = \operatorname{sign}(W^2)\sqrt{|W^2|},\qquad
\cos\theta_{pq} = \hat p_h\cdot\hat q,\nonumber\\
t &= Q^2-M_h^2+2\nu E_h-2P_h|\vec q|\cos\theta_{pq},\qquad t_\mathrm{min} = Q^2-M_h^2+2\nu E_h-2P_h|\vec q|,
\lbl{eq:mesonkin}
\end{feqa}
$\varepsilon$ of \ceqr{eq:epsilon} and $\varphi_{pq}$ of \ceqr{eq:phipq} ($t$ is $-t_\mathrm{Mandelstam}$, positive). For
$\pi^0$ the two decay photons are drawn here (two random numbers; the photon tracking is described with the pion
reaction). The missing quantities are those of the whole nucleus:
\begin{feqa}{reactions/sidis/SidisProduction.cpp: SidisProduction::completeEv}{event.f:914-918, 952-954}
\vec P_m &= P_h\hat p_h-|\vec q|\,\hat q,\qquad P_m = P_\mathrm{miss}=|\vec P_m|,\qquad E_m = E_\mathrm{miss} = \nu+M_A-E_h,
\lbl{eq:pmiss}
\end{feqa}
so the vertex ``missing energy'' of SIDIS is the energy of the unobserved system, not a separation energy; the residual
mass and kinetic energy are not computed. The try fails below a two-pion threshold evaluated on the struck nucleon at rest:
\begin{feq}{reactions/sidis/SidisProduction.cpp: SidisProduction::completeEv}{event.f:979-988}
(M_N+\nu-E_h)^2 - P_m^2 < (M_p+M_{\pi^0})^2\ \Rightarrow\ \text{fail},
\lbl{eq:threshold}
\end{feq}
with $M_{\pi^0}=134.9766$ MeV for every hadron (also for kaons, whose lowest semi-inclusive final state is a hyperon),
and without the Fermi momentum. Then
\begin{feq}{reactions/sidis/SidisProduction.cpp: SidisProduction::completeEv}{event.f:990-998}
z = E_h/\nu,\qquad p_T^2 = P_h^2\bigl(1-\cos^2\theta_{pq}\bigr),\qquad z>1\ \Rightarrow\ \text{fail (with a warning)},
\lbl{eq:zpt}
\end{feq}
the last test being unreachable because \eqr{eq:threshold} already rejects $z>1$; the warning line of the code is
therefore never printed by a run. The vertex $P_{m,\perp}$ of \ceqr{eq:pmcomp} is computed before $P_m$ is assigned, i.e.\
from the previous try's $P_m$ in the Fortran; SIMC++ uses the current $P_m$ by default (it is not a tree column).

\subsection{Reconstruction}
\begin{feqa}{reactions/EepReactionBase.cpp: EepReactionBase::completeRecon}{event.f:1325-1341}
E_m^\mathrm{rec} &= \nu^\mathrm{rec}+M_N-E_h^\mathrm{rec},\qquad
M_X^\mathrm{rec} = \operatorname{sign}(m^2)\sqrt{|m^2|},\quad m^2 = (E_m^\mathrm{rec})^2-(P_m^\mathrm{rec})^2,\nonumber\\
M_{X,A}^\mathrm{rec} &= \operatorname{sign}(m_A^2)\sqrt{|m_A^2|},\quad m_A^2 = (\nu^\mathrm{rec}+M_A-E_h^\mathrm{rec})^2-(P_m^\mathrm{rec})^2,\nonumber\\
t^\mathrm{rec} &= (Q^\mathrm{rec})^2-M_h^2+2\bigl(\nu^\mathrm{rec}E_h^\mathrm{rec}-P_h^\mathrm{rec}|\vec q^\mathrm{rec}|\cos\theta^\mathrm{rec}_{pq}\bigr),\nonumber\\
z^\mathrm{rec} &= E_h^\mathrm{rec}/\nu^\mathrm{rec},\qquad p_T^{2,\mathrm{rec}} = (P_h^\mathrm{rec})^2\bigl(1-\cos^2\theta^\mathrm{rec}_{pq}\bigr),
\lbl{eq:recon}
\end{feqa}
$M_X^\mathrm{rec}$ (tree column \code{missmass}) being the missing mass on a nucleon at rest and $M_{X,A}^\mathrm{rec}$
(\code{mmnuc}) that on the nucleus. $W^\mathrm{rec}$ and $x^\mathrm{rec}$ use the proton mass (\ceqr{eq:w}).

% =====================================================================================================================
\section{Jacobians}\lbl{sec:jac}
Both pairs of slopes are generated flat, so the event jacobian is the product of the two solid-angle factors of
\ceqr{eq:jacomega}:
\begin{feq}{reactions/sidis/SidisProduction.cpp: SidisProduction::completeEv}{event.f:1015-1025}
J = \bigl(1+x_e'^2+y_e'^2\bigr)^{-3/2}\bigl(1+x_h'^2+y_h'^2\bigr)^{-3/2}.
\lbl{eq:jac}
\end{feq}
The energies are generated flat in $E'$ and $E_h$, the variables of the cross section; the change of variables from the
model's natural $(z,p_T^2)$ to $(E_h,\cos\theta_{pq})$ is inside the model (\eqr{eq:xsjac}).

% =====================================================================================================================
\section{Cross-section model}\lbl{sec:xsec}
\subsection{What the reaction requires from a model}
\code{SidisProduction::completeMain} sets $S=1$ and $n_\mathrm{tgt}=1$ (the sum over the nucleons is inside the model)
and calls \code{sigccMain(vertex, main)}, which may write event fields (the tree columns \code{sighad} and
\code{jacobian}, the survival probability). The model reads $E_\mathrm{in}$, $E'$, $\nu$, $Q^2$, $\vec q$, $E_h$, $P_h$,
$\theta_{pq}$, $z$, $p_T^2$ of the vertex record and, with Fermi motion, the Fermi momentum of the try; it must return
$\mathrm{d}^5\sigma/(\mathrm{d}\Omega_e\,\mathrm{d}E'\,\mathrm{d}\Omega_h\,\mathrm{d}E_h)$ per nucleus in
$\mu$b/(MeV$^2$\,sr$^2$), the units of \eqr{eq:genvol}. The weight is
\begin{feq}{reactions/sidis/SidisProduction.cpp: SidisProduction::completeMain}{event.f:1561-1564}
w = J\cdot w_\mathrm{gen}\cdot\sigma\cdot\bigl(1+V_C/E_b\bigr)^2_{[\text{with \code{using\_Coulomb}}]}\ \cdot P_\mathrm{surv}{}_{[\text{K}^\pm,\ \text{decay not simulated}]}
\quad[\mu\mathrm{b/(MeV^2\,sr^2)}],
\lbl{eq:weight}
\end{feq}
with $P_\mathrm{surv}$ of \eqr{eq:survival}; the pion survival probability is computed but not applied.

\subsection{The model \code{sidis}: overview}
\code{peepiX} works in GeV. The code comment states the idea: ``p(e,e\pr$\pi^+$)X semi-inclusive cross sections from the
CTEQ5M parton distribution functions and a simple parameterization of the favored and unfavored fragmentation
functions'', later replaced by the fits below. The cross section is
\begin{feq}{physics/sidis/SidisCrossSection.cpp: peepiX}{semi\_physics.f:508-513, 525-526, 553-556}
\frac{\mathrm{d}^5\sigma}{\mathrm{d}\Omega_e\,\mathrm{d}E'\,\mathrm{d}z\,\mathrm{d}p_T^2\,\mathrm{d}\varphi_{pq}}
= \sigma_e\cdot R_h(A,z)\,\frac{1}{z}\frac{\sum_q e_q^2\,q_A(x)\,zD_q^h(z)}{\sum_q e_q^2\,q_A(x)}\cdot
\frac{b\,e^{-b\,p_T^2}}{2\pi}\cdot(\hbar c)^2\cdot10^4\quad[\mu\mathrm{b/(GeV^3\,sr)}],
\lbl{eq:sigsemi}
\end{feq}
with the sums over $q=u,\bar u,d,\bar d,s,\bar s$, $e_u=2/3$, $e_d=e_s=-1/3$, the inclusive cross section $\sigma_e$
(\eqr{eq:sige}) and the $p_T$ slope
\begin{feq}{physics/sidis/SidisCrossSection.cpp: peepiX}{semi\_physics.f:525}
b = \frac{1}{0.120\,z^2+0.200}\quad[\mathrm{GeV^{-2}}],
\lbl{eq:bslope}
\end{feq}
which normalises the $p_T^2$ distribution to 1, so that $\int\mathrm{d}p_T^2\,\mathrm{d}\varphi_{pq}$ of the middle factors is
the multiplicity $R_h\,\mathrm{d}N^h/\mathrm{d}z$. The factor $(\hbar c)^2\cdot10^4$ converts GeV$^{-2}$ to $\mu$b
($\hbar c$ in GeV\,fm). The tree column \code{sighad} is $R_h\,(\ldots)\,b\,e^{-bp_T^2}/2\pi$ in GeV$^{-2}$, i.e.\
\eqr{eq:sigsemi} without $\sigma_e$ and the conversion factor.

\subsection{Kinematics of the model}\lbl{sec:fermixs}
\begin{feqa}{physics/sidis/SidisCrossSection.cpp: peepiX}{semi\_physics.f:237-258, 268-272}
E_h &= z\nu,\qquad P_h = \sqrt{E_h^2-M_h^2},\qquad \cos\theta_{pq}\ \text{from the vertex},\nonumber\\
x &= \frac{Q^2}{2M_N\nu}\quad\text{or, with Fermi motion,}\quad x = \frac{Q^2}{2\,(E_F\nu-p_F\,\hat p_F\cdot\vec q)},\qquad
x\leftarrow\min(x,1),\nonumber\\
W^2 &= M_p^2+Q^2\Bigl(\frac1x-1\Bigr),
\lbl{eq:xsidis}
\end{feqa}
with $M_N=M_p$; $x>1$ prints ``XBj is too large!'' and is replaced by 1. $M_h$ is the charged-pion mass also for a $\pi^0$
in the Fortran (corrected to $M_{\pi^0}$ in the default mode). With Fermi motion ($p_F$, $E_F$ of \eqr{eq:fermi},
deuterium only) $x$ is Bjorken's variable of the moving nucleon, $Q^2/2P\cdot q$; this $x$ is used everywhere in the model
and $W^2$ is computed from it. A non-positive $x$ (a large Fermi momentum along $\vec q$) stops the Fortran run in the
parton-distribution routine; SIMC++ returns $\sigma=0$ for that event in the default mode. A second missing-mass cut
returns $\sigma=0$:
\begin{feq}{physics/sidis/SidisCrossSection.cpp: peepiX}{semi\_physics.f:276-284}
M_N^2+2M_N\nu(1-z)\bigl(1-p_T^2[\mathrm{GeV^2}]\bigr) < (M_N+M_h)^2\ \Rightarrow\ \sigma = 0\qquad(\text{masses and }\nu\text{ in GeV}),
\lbl{eq:mmcut}
\end{feq}
described by the code as the ``2$\pi$ / $\pi$K threshold''. The left-hand side is the missing mass squared of the
collinear limit, $M_N^2+2M_N\nu(1-z)$, multiplied in its second term by the dimension-carrying factor $1-p_T^2$; it is
not a missing mass for $p_T\ne0$ (the code comment adds ``this shouldn't be necessary''). Below the cut the model returns
before it writes \code{sighad}, the jacobian column and the survival probability: \code{sighad} stays 0 (reset at every
try), the jacobian column keeps the previous event's value, and the weight is 0.

\subsection{Parton distributions: CTEQ5M}
The model reads the CTEQ5M table \citer{cteq5} (\code{cteq5m.tbl}, the set the code calls \code{iset} = 1; MS-bar, NLO,
$\Lambda^{(5)}_\mathrm{QCD}=0.226$ GeV) and calls it with $Q=\sqrt{Q^2[\mathrm{GeV^2}]}$ for $u,\bar u,d,\bar d,s,\bar s$ of
the proton. The table holds $f(x,Q)$ on a grid of 91 $x$ nodes $x_j$ and 14 nodes of $\ln(Q/\Lambda)$ for the flavours
$\bar b\ldots\bar u$ (sea, quark $=$ antiquark for $s,c,b$), $g$, $u$, $d$; the value is a two-dimensional quadratic
(three-point Neville) interpolation, first in $x$ at three $Q$ nodes, then in $\ln(Q/\Lambda)$:
\begin{feqa}{physics/sidis/Cteq5Pdf.cpp: Cteq5Pdf::partonX, polint}{cteq5/Ctq5Pdf.f:133-198, 313-351}
j_x &= j_\mathrm{lo}(x),\ \ \text{clamped to }[0,\,N_x-2],\qquad j_q = j_\mathrm{lo}(\ln(Q/\Lambda)),\ \ \text{clamped to }[0,\,N_t-2],\nonumber\\
F_i &= \mathcal{P}_2\bigl(x;\ \{x_{j_x+k},\ f(x_{j_x+k},Q_{j_q+i})\}_{k=0,1,2}\bigr),\quad i=0,1,2,\qquad
f(x,Q) = \mathcal{P}_2\bigl(\ln\tfrac{Q}{\Lambda};\ \{\ln\tfrac{Q_{j_q+i}}{\Lambda},\ F_i\}\bigr),
\lbl{eq:cteq}
\end{feqa}
with $j_\mathrm{lo}$ the index of the lower node of the interval containing the argument (bisection) and $\mathcal P_2$
the quadratic through the three points; $s$ is evaluated as $\bar s$. A negative result is set to 0
\fsrc{cteq5/Ctq5Pdf.f:102}. $x<0$ or $x>1$ stops the run; $Q<\Lambda$ prints a warning and replaces the caller's $Q$ by
$\Lambda$ for the remaining calls of the event. Below the smallest node $x_\mathrm{min}$ the clamp of $j_x$ is applied
only in the first such call of a run (with a warning); the tables were not compared with the CTEQ publication for this
note.

\begin{simctable}{footnotesize}{L{0.35\linewidth}L{0.65\linewidth}}
\caption{The CTEQ5M table as read (\code{ReadTbl}; one property per row).}\lbl{tab:cteq}\\
\toprule\rowcolor{simchead}\simcth{property} & \simcth{value}\\\midrule\endfirsthead
\bottomrule\endlastfoot
order, flavours & 2 (NLO), 5\\
$\Lambda$ & 0.2260 GeV\\
quark masses $m_c$, $m_b$, $m_t$ & 1.3, 4.5, 180 GeV\\
$N_x$, $N_t$, flavours stored & 90, 13 (indices from 0), $N_f+3 = 8$\\
$Q$ range of the grid & 1 to $10^4$ GeV (extrapolated outside)\\
\end{simctable}

The nucleus is an incoherent sum of $Z$ protons and $N=A-Z$ neutrons with isospin symmetry:
\begin{feqa}{physics/sidis/SidisCrossSection.cpp: peepiX}{semi\_physics.f:322-329}
u_A &= Zu+Nd,\quad \bar u_A = Z\bar u+N\bar d,\quad d_A = Zd+Nu,\quad \bar d_A = Z\bar d+N\bar u,\nonumber\\
s_A &= (Z+N)s,\quad \bar s_A = (Z+N)\bar s,\nonumber\\
\Sigma_{e^2} &= e_u^2(u_A+\bar u_A)+e_d^2(d_A+\bar d_A)+e_s^2(s_A+\bar s_A).
\lbl{eq:pdfA}
\end{feqa}
When every distribution vanishes ($x=1$ after the clamp) the ratio of \eqr{eq:sigsemi} is $0/0$: the Fortran weight is
NaN, SIMC++ uses 0 in the default mode.

\subsection{Pion fragmentation: P.\ Bosted's fit}
For pions the code uses ``Peter Bosted's parameterization'' (fit of 9/20/2021 to Hall C data, a modified scaling
variable after Accardi et al.\ \citer{accardi}); the fit gives $zD$ directly. With $M=M_p$, $m=M_\pi$ (GeV) and $W$ from
\eqr{eq:xsidis}:
\begin{feqa}{physics/sidis/SidisCrossSection.cpp: peepiX}{semi\_physics.f:465-477}
\xi &= \frac{2x}{1+\sqrt{1+4x^2M^2/Q^2}},\qquad s_v = \ln(Q^2/2),\nonumber\\
z_p &= \frac z2\,\frac{\xi}{x}\Bigl(1+\sqrt{1-\frac{4x^2M^2(m^2+p_T^2)}{z^2Q^4}}\Bigr),\nonumber\\
y_k &= c_{k,1}\,z_p^{\,c_{k,2}+c_{k,4}s_v+c_{k,9}/W}\,(1-z_p)^{c_{k,3}+c_{k,5}s_v+c_{k,10}/W}\nonumber\\
&\quad\times\bigl(1+c_{k,6}z_p+c_{k,7}z_p^2+c_{k,8}z_p^3\bigr)\Bigl(1+\frac{c_{k,11}}{W}+\frac{c_{k,12}}{W^2}\Bigr),
\lbl{eq:bosted}
\end{feqa}
$k=f$ (favoured) or $u$ (unfavoured), with the coefficients of \tabr{tab:pbfit}. $\xi$ is the Nachtmann variable. With
$x=Q^2/2M\nu$ the argument of the inner root equals $1-(m^2+p_T^2)/E_h^2 = p_L^2/E_h^2$ ($p_L$ the hadron momentum along
$\vec q$), so $z_p=\tfrac z2\tfrac\xi x(1+p_L/E_h)$ is the light-cone fraction corrected for the target mass; with the
Fermi $x$ or the clamp $x=1$ this identity no longer holds. The assignment to the quark flavours is
\begin{feqa}{physics/sidis/SidisCrossSection.cpp: peepiX}{semi\_physics.f:479-495}
\pi^+:&\quad zD_u = zD_{\bar d} = y_f,\qquad zD_d = zD_{\bar u} = zD_s = zD_{\bar s} = y_u,\nonumber\\
\pi^-:&\quad zD_d = zD_{\bar u} = y_f,\qquad zD_u = zD_{\bar d} = zD_s = zD_{\bar s} = y_u,\nonumber\\
\pi^0:&\quad zD_u = zD_d = zD_{\bar u} = \tfrac12(y_f+y_u),\qquad zD_{\bar d} = \text{the }\pi^\pm\text{ value above},\qquad zD_s=zD_{\bar s}=y_u.
\lbl{eq:pionff}
\end{feqa}
The $\pi^0$ block is introduced by the comment ``just get averge of pi+ and pi-'' but averages $D_d$ twice and never
$D_{\bar d}$, which keeps the value of the charged pion selected by \code{doing\_hplus}; SIMC++ averages $D_{\bar d}$ too
in the default mode. The coefficients are transcribed from the code; they were not compared with a publication (the fit
is not published in the code).

\begin{simctable}{footnotesize}{L{0.1\linewidth}L{0.2\linewidth}L{0.2\linewidth}L{0.5\linewidth}}
\caption{Coefficients of \eqr{eq:bosted} (\fsrc{semi\_physics.f:107-112}, ``PB fit of 9/20/2021''; one coefficient per row).}\lbl{tab:pbfit}\\
\toprule\rowcolor{simchead}\simcth{$i$} & \simcth{$c_{f,i}$} & \simcth{$c_{u,i}$} & \simcth{role}\\\midrule\endfirsthead
\toprule\rowcolor{simchead}\simcth{$i$} & \simcth{$c_{f,i}$} & \simcth{$c_{u,i}$} & \simcth{role}\\\midrule\endhead
\bottomrule\endlastfoot
1 & 1.0424 & 0.7840 & normalisation\\
2 & $-0.1714$ & 0.2369 & power of $z_p$\\
3 & 1.8960 & 1.4238 & power of $1-z_p$\\
4 & $-0.0307$ & 0.1484 & $s_v$ term of the $z_p$ power\\
5 & 0.1636 & 0.1518 & $s_v$ term of the $1-z_p$ power\\
6 & $-0.1272$ & $-1.2923$ & linear term of the polynomial\\
7 & $-4.2093$ & $-1.5710$ & quadratic term\\
8 & 5.0103 & 3.0305 & cubic term\\
9 & 2.7406 & 1.1995 & $1/W$ term of the $z_p$ power\\
10 & $-0.5778$ & 1.3553 & $1/W$ term of the $1-z_p$ power\\
11 & 3.5292 & 2.5868 & $1/W$ factor\\
12 & 7.3910 & 8.0666 & $1/W^2$ factor\\
\end{simctable}

\subsection{Kaon fragmentation: fDSS}
For kaons the model calls the fDSS routine of de Florian, Sassot and Stratmann \citer{dss} at NLO for $K^\pm$
(\code{KANLO.GRID}, read at start-up), with $z$ and $Q^2$; it returns $z$ times the fragmentation functions. The grid
holds, at 34 $z$ nodes and 24 $Q^2$ nodes, the totals $zD_{q+\bar q}$ for $u$, $d$, $s$, $c$, $b$, the gluon, and the
valence combinations $zD_{q-\bar q}$ for $u$, $d$, $s$. Each is divided at start-up by a smooth factor and the quotient is
interpolated bilinearly in $(\ln z,\ln Q^2)$ (CERNLIB E104, linear extrapolation outside the grid, a node at $z=1$ set
to 0):
\begin{feqa}{physics/sidis/Fdss.cpp: FdssGrid::load, FdssGrid::eval, ffint}{fdss/fdss.f:117-161, 213-257}
\tilde T(z_j,Q^2_k) &= \frac{T(z_j,Q^2_k)}{(1-z_j)^{a}\,z_j^{\,c}},\qquad
T(z,Q^2) = \mathcal{I}\bigl[\tilde T\bigr](\ln z,\ln Q^2)\cdot(1-z)^{a}z^{c},\nonumber\\
(a,c) &= (4,\,0.5)\ \text{for }u,d,s\text{ (total and valence)},\quad (7,\,0.3)\ \text{for }c,b,\quad (4,\,0.3)\ \text{for }g,\nonumber\\
zD_q &= \tfrac12(T_q+V_q),\qquad zD_{\bar q} = \tfrac12(T_q-V_q),\qquad zD_c = \tfrac12T_c,\qquad zD_b=\tfrac12T_b,
\lbl{eq:fdss}
\end{feqa}
with $\mathcal I$ the bilinear interpolation, $T_q$ the total and $V_q$ the valence combination. These are the $K^+$
functions; for $K^-$ quark and antiquark are exchanged. A $z$ outside $[0.05,1]$ or a $Q^2$ outside $[1,10^5]$ GeV$^2$
prints a warning and is extrapolated. The Fortran evaluates the factors $(1-z)^n$ and the products with the
interpolated value in quadruple precision (its double-precision literals are promoted by the build); SIMC++ reproduces
that arithmetic. The grid is the published one; its values were not compared with the paper.

The partonic sum of \eqr{eq:sigsemi} is
\begin{feq}{physics/sidis/SidisCrossSection.cpp: peepiX}{semi\_physics.f:508-513}
\frac{\mathrm dN}{\mathrm dz} = \frac{e_u^2\,(u_A\,zD_u+\bar u_A\,zD_{\bar u}) + e_d^2\,(d_A\,zD_d+\bar d_A\,zD_{\bar d})
+ e_s^2\,(s_A\,zD_s+\bar s_A\,zD_{\bar s})}{\Sigma_{e^2}\;z}.
\lbl{eq:dndz}
\end{feq}

\subsection{Nuclear factors}
The multiplicity is multiplied by a hadron-attenuation fit that is 1 only for $A=1$ (deuterium gets
$R_h\approx0.955$ at $z=0.5$):
\begin{feqa}{physics/sidis/SidisCrossSection.cpp: rhadGlobal}{semi\_physics.f:652-680}
R_h(A,z) &= N_0\,z^{\alpha_h}(1-z)^{\beta_h}\quad(A\ne1),\qquad R_h(1,z)=1,\nonumber\\
N_0 &= 0.98883-0.0038309A+1.0841\times10^{-5}A^2,\nonumber\\
\alpha_h &= 0.031953-0.0018659A+5.1747\times10^{-6}A^2,\nonumber\\
\beta_h &= 0.0085475+0.0012763\tilde A-2.4451\times10^{-6}\tilde A^2,\qquad \tilde A=\min(A,83.8).
\lbl{eq:rhad}
\end{feqa}
The inclusive part (below) carries the SLAC EMC ratio for $A>2$ \citer{gomez}:
\begin{feqa}{physics/sidis/SidisCrossSection.cpp: slacemcfit}{semi\_physics.f:634-650}
R_\mathrm{EMC}(A,x) &= C(x)\,A^{\alpha(x)}\ \ (A>2),\qquad R_\mathrm{EMC}=1\ \ (A\le2),\nonumber\\
\alpha(x) &= -0.070+2.189x-24.667x^2+145.291x^3-497.237x^4\nonumber\\
&\quad+1013.129x^5-1208.393x^6+775.767x^7-205.872x^8,\nonumber\\
C(x) &= \exp\bigl(0.017+0.018\ln x+0.005\ln^2x\bigr).
\lbl{eq:emc}
\end{feqa}
Neither fit names its source in the code; the coefficients are transcribed from the code and not compared with a paper.

\subsection{Inclusive cross section}
The inclusive cross section is the one-photon formula with the structure functions of the Christy--Bosted-type fit
\code{F1F2IN21} (E.\ Christy, 2021) \citer{christy}, evaluated per nucleon at $W^2$ of \eqr{eq:xsidis} and summed:
\begin{feqa}{physics/sidis/SidisCrossSection.cpp: peepiX}{semi\_physics.f:528-553}
F_{1,2} &= \bigl(Z\,F_{1,2}^p+N\,F_{1,2}^n\bigr)\,R_\mathrm{EMC}(A,x),\qquad W_1 = F_1/M_N,\qquad W_2 = F_2/\nu,\nonumber\\
\sin^2\frac{\theta_e}{2} &= \frac{Q^2}{4E_\mathrm{in}E'},\qquad
\sigma_e = \frac{4\alpha^2E'^2}{Q^4}\Bigl[W_2\cos^2\frac{\theta_e}{2}+2W_1\sin^2\frac{\theta_e}{2}\Bigr]\quad[\mathrm{GeV^{-3}\,sr^{-1}}],
\lbl{eq:sige}
\end{feqa}
in GeV ($\sigma_e$ is $\mathrm{d}^2\sigma/\mathrm{d}\Omega_e\mathrm{d}E'$ per nucleus in natural units). The free-nucleon
structure functions come from fits of the transverse and longitudinal virtual-photon cross sections:
\begin{feqa}{physics/sidis/F1F2In21.cpp: sf}{F1F2IN21\_v1.0.f:2345-2371}
x_W &= \frac{Q^2}{Q^2+W^2-M^2},\qquad
F_1 = \frac{\sigma_T\,|W^2-M^2|}{0.3894\times10^3\cdot8\pi^2\alpha},\qquad
F_L = \frac{2x_W\,\sigma_L\,|W^2-M^2|}{0.3894\times10^3\cdot8\pi^2\alpha},\nonumber\\
F_2 &= \frac{2x_WF_1+F_L}{1+4M^2x_W^2/Q^2},
\lbl{eq:sfconv}
\end{feqa}
($M=0.938272$ GeV, $\pi=3.14159$, $\alpha=1/137.03599$ in this routine; $\sigma_{T,L}$ in $\mu$b). $\sigma_{T,L}$ of the proton
and the neutron are the sum of seven Breit--Wigner resonances and a non-resonant term:
\begin{feqa}{physics/sidis/F1F2In21.cpp: resmod}{F1F2IN21\_v1.0.f:373-670, 672-959}
\sigma &= W\sum_{i=1}^{7} h_i\,\frac{k_r}{k}\frac{k^\mathrm{cm}_r}{k^\mathrm{cm}}\,
\frac{\Gamma_i\,\Gamma_{\gamma,i}/\Gamma_i^0}{(W^2-M_i^2)^2+M_i^2\Gamma_i^2}\ [\times Q^2\text{ for }\sigma_L]
+389.4\,A_0\,(1-x_\pi)^{t_1}x_\pi^{t_2}\,[\ldots],\nonumber\\
x_\pi &= \Bigl(1+\frac{W^2-(M+m_{\pi^0})^2}{Q^2+Q_0^2}\Bigr)^{-1},
\lbl{eq:resmod}
\end{feqa}
with energy-dependent widths $\Gamma_i$ built from the one-pion, two-pion and $\eta$ decay momenta (branching ratios and
angular momenta of P33(1232), S11(1535), D13(1520), F15(1680), S11(1650), P11(1440), F37(1950) fixed in the code), $Q^2$
dependent heights $h_i$, and an $\eta$-threshold term in $\sigma_T$. The 100 parameters of each nucleon (masses, widths,
heights, non-resonant coefficients) are DATA statements of \code{rescsp} and \code{rescsn}; they are not reproduced here and
were not compared with a publication. For $A=1$ only the free-nucleon branch of \code{F1F2IN21} is used (the deuteron and
nuclear branches of the routine are never called by SIDIS); its first call also reads deuteron smearing tables, which have
no effect here and are not read by SIMC++.

\subsection{Change of variables, Fermi factor and units}
The event is generated in $(E_h,\Omega_h)$; with $\Omega_h$ measured around $\vec q$ ($\mathrm{d}\Omega_h =
\mathrm{d}\cos\theta_{pq}\,\mathrm{d}\varphi_{pq}$), $z=E_h/\nu$ and $p_T^2 = P_h^2(1-\cos^2\theta_{pq})$:
\begin{feqa}{physics/sidis/SidisCrossSection.cpp: peepiX}{semi\_physics.f:562-586}
J_{zp_T} &= \Bigl|\frac{\partial(z,p_T^2)}{\partial(E_h,\cos\theta_{pq})}\Bigr| = \frac{2P_h^2\cos\theta_{pq}}{\nu}\quad[\mathrm{GeV}],\nonumber\\
\sigma &= \frac{\mathrm{d}^5\sigma}{\mathrm{d}\Omega_e\,\mathrm{d}E'\,\mathrm{d}z\,\mathrm{d}p_T^2\,\mathrm{d}\varphi_{pq}}\cdot J_{zp_T}\cdot10^{-6}\cdot f_F
\quad[\mu\mathrm{b/(MeV^2\,sr^2)}],\nonumber\\
f_F &= \frac{M_N}{E_F}\,\frac{1}{1-p_F\hat p_{F,z}/E_F}\ \ (\text{Fermi motion}),\qquad f_F=1\ \ (\text{otherwise}),
\lbl{eq:xsjac}
\end{feqa}
$10^{-6}$ converting GeV$^{-2}$ to MeV$^{-2}$. The code writes $\cos\theta_{pq}$, not its absolute value; it is positive in
any spectrometer acceptance near $\vec q$. $f_F$ is the flux factor of a non-collinear collision, explained in the code:
the relative velocity becomes $1-p_{F,z}/E_F$ and the target energy $E_F$ instead of $M_N$. The tree column
\code{jacobian} is $1000\,J_{zp_T}$ (MeV). The header comment of the routine gives the units of an older version
(``d3sigma/dEe'dOmegae'Omegapi, microbarn/MeV/sr\^{}2''); the returned value has the units above.

\subsection{Survival probability}
When \code{doing\_decay} = 0 the model computes the probability that the hadron reaches the aerogel detector of the
hadron arm, written to the tree column \code{decdist}; it multiplies the weight of kaons only (\eqr{eq:weight}):
\begin{feqa}{physics/sidis/SidisCrossSection.cpp: peepiX}{semi\_physics.f:610-626}
L &= L_\mathrm{FP} + z_\mathrm{aero}\bigl(1+x'^2_\mathrm{fp}+y'^2_\mathrm{fp}\bigr),\qquad
\beta = \frac{P}{\sqrt{P^2+M_h^2}},\qquad \gamma = (1-\beta^2)^{-1/2},\nonumber\\
P_\mathrm{surv} &= \exp\bigl(-L/(c\tau\,\beta\gamma)\bigr),
\lbl{eq:survival}
\end{feqa}
with $L_\mathrm{FP}$ the path length of the hadron to the end of the hut, $x'_\mathrm{fp}$, $y'_\mathrm{fp}$ the focal-plane
slopes, $c\tau$ the input \code{ctau}, and $z_\mathrm{aero}$ of \tabr{tab:zaero}. The Fortran uses the central momentum
$P=P_h^0$ although $P_h^0(1+\delta^\mathrm{SP}_h/100)$ is computed next to it (the comment says ``Get particle momentum from
main.SP.p.delta''), and $z_\mathrm{aero}(1+x'^2+y'^2)$ where the path along a slope is $z_\mathrm{aero}\sqrt{1+x'^2+y'^2}$;
SIMC++ uses the hadron momentum and the square root in the default mode.

\begin{simctable}{footnotesize}{L{0.25\linewidth}L{0.2\linewidth}L{0.55\linewidth}}
\caption{Aerogel distance $z_\mathrm{aero}$ of \eqr{eq:survival} (one arm per row).}\lbl{tab:zaero}\\
\toprule\rowcolor{simchead}\simcth{hadron arm} & \simcth{$z_\mathrm{aero}$ [cm]} & \simcth{comment in the code}\\\midrule\endfirsthead
\bottomrule\endlastfoot
1 (HMS) & $-331.491$ & ``aero at 40.199 cm, last project=371.69 cm''\\
2 (SOS) & $-82.8$ & ``From Rick: aero at 263.2cm, last project=346(cal)''\\
3 (HRS-R) & $-183.0$ & ``aero at 130cm, last project=313(S2)''\\
4 (HRS-L) & $-183.0$ & as HRS-R\\
5 (SHMS) & none & the previous value is kept (0 in a run); 0 in the SIMC++ default mode\\
\end{simctable}

\subsection{Central cross section and unused parts}
The \texttt{.hist} line \code{CENTRAL.sigcc} is the model at the central event, completed by the reconstruction formulas
(so $z$, $p_T^2$ and $\theta_{pq}$ of the central ray are defined), and the run log prints its kinematics and the
multiplicity \code{sighad} (``SIDIS: central event kinematics''); the ``lab cross section (nb/GeV2/sr2)'' printed there is
$10^9\,\sigma$. The model has a second, ``central'' branch that evaluates \eqr{eq:sigsemi} at a fixed $z$ or $p_T^2$ in
bins (inputs \code{sigc\_flag}, \code{sigc\_nbin}, \code{sigc\_kin\_min/max/ind}); its call is commented out, so these
inputs have no effect (SIMC++ ports the branch, never calls it, and warns when they are set); the tree columns
\code{sigcent} and \code{centjac} of the legacy layout are always 0. The input \code{pt\_b\_param} is read and not used
(\eqr{eq:bslope} replaces it). The dilution-factor routine is not called either; it recomputes a $\pi^\pm$ multiplicity
with the older Binnewies parametrisation of the fragmentation functions \citer{binnewies} and returns
\begin{feqa}{physics/sidis/SemiDilution.cpp: semiDilution}{semi\_dilution.f:239-251}
f_\mathrm{dil} &= \frac{\sigma_\mathrm{tgt}}{\sigma_\mathrm{tgt}+0.3333\,\sigma_{^{15}\mathrm N}+0.3165\,\sigma_{^4\mathrm{He}}
+0.0144\,\sigma_\mathrm{Al}+0.0031\,\sigma_\mathrm{Cu}+0.0013\,\sigma_\mathrm{Ni}},\nonumber\\
\sigma_{(Z,N)} &= Z\sigma_p+N\sigma_n,
\lbl{eq:dilution}
\end{feqa}
(999 if the proton or neutron term vanishes), the dilution of a polarised $^{15}$NH$_3$-type target; SIMC++ keeps it as a
tested function only.

% =====================================================================================================================
\section{Radiation, energy loss and decays}\lbl{sec:rad}
The three tails radiate as in \csecr{sec:radiation}; the hadron tail uses $\lambda_3$ with the hadron's $E_h$, $P_h$
(\ceqr{eq:lambda}). The photon limits are those of the meson reactions:
\begin{feqa}{radiation/Radc.cpp: generateRad}{radc.f:297-308, 372-389, 428-446}
E_{\gamma1} &\in \bigl[-d_E,\ \min(\Sigma_\mathrm{max}-E',\ E^\mathrm{max}_{\gamma1})+d_E\bigr],\nonumber\\
E_{\gamma2} &\in \bigl[E'-E'_\mathrm{max}-d_E,\ \min(E'-E'_\mathrm{min},\ E_\gamma^\mathrm{tot,max}-E_{\gamma1})+d_E\bigr],\nonumber\\
E_{\gamma3} &\in \bigl[E_h-E_{h,\mathrm{max}}-d_E,\ \min(E_h-E_{h,\mathrm{min}},\ E_\gamma^\mathrm{tot,max}-E_{\gamma1}-E_{\gamma2})+d_E\bigr],
\lbl{eq:radlimits}
\end{feqa}
(a negative lower limit counts as 0). The missing-energy constraints of the (e,e\pr p) reactions do not apply. After a
tail-1 photon \code{completeEv} is called again at the lower beam energy: the thresholds of \eqsr{eq:threshold}{eq:zpt}
are re-applied, the generated $E'$, $E_h$ and directions are kept. The soft-photon factors use elastic e--p kinematics
(\ceqr{eq:bremkin}), an approximation for this inelastic reaction.

The hadron loses energy and scatters in the target and the windows with its own mass (\csecr{sec:eloss}); the Coulomb
correction applies to the electron for $Z>1$ (\csecr{sec:coulomb}). With \code{doing\_decay} = 1 a charged pion or kaon may
decay in flight inside the hadron arm (the spectrometer models follow the decay product; the column \code{decdist} then
holds the decay distance and \code{Mhadron} the mass of the particle that reached the detectors); with
\code{doing\_decay} = 0 the survival probability of \eqr{eq:survival} is used. A $\pi^0$ decays at the vertex into two
photons (\code{pizero\_ngamma} = 1 or 2 photons required in the calorimeter arm; the photon columns are added to the tree).

% =====================================================================================================================
\section{Inputs and outputs}\lbl{sec:io}
\subsection{Inputs}
\begin{simctable}{footnotesize}{L{0.24\linewidth}L{0.12\linewidth}L{0.08\linewidth}L{0.56\linewidth}}
\caption{Inputs specific to SIDIS (one input per row; the common inputs are in note 01, \S7, and the user manual,
chapter 4).}\lbl{tab:inputs}\\
\toprule\rowcolor{simchead}\simcth{input} & \simcth{unit} & \simcth{default} & \simcth{meaning here}\\\midrule\endfirsthead
\toprule\rowcolor{simchead}\simcth{input} & \simcth{unit} & \simcth{default} & \simcth{meaning here}\\\midrule\endhead
\bottomrule\endlastfoot
\code{doing\_semi} & flag & 0 & 1 selects this reaction\\
\code{doing\_pion} & flag & 0 & 1: the hadron is a pion\\
\code{doing\_kaon} & flag & 0 & 1: the hadron is a kaon\\
\code{doing\_hplus} & flag & 0 & 1: positive hadron, 0: negative\\
\code{doing\_pizero} & flag & 0 & 1 (with \code{doing\_pion}): $\pi^0$\\
\code{pizero\_ngamma} & -- & 0 & $\pi^0$: number of photons required (1 or 2)\\
\code{targ\%A} & -- & 0 & 1 H, 2 D, $\ge3$ nucleus; $A$ of the nuclear factors\\
\code{targ\%Z} & -- & 0 & $Z$ of the nucleon sum (\eqr{eq:pdfA}); Coulomb correction for $Z>1$\\
\code{do\_fermi} & flag & 0 & deuterium: Fermi motion in the cross section (\secr{sec:fermixs}); ignored for $A\ge3$, refused (switched off) for H\\
\code{doing\_decay} & flag & 0 & 1: decay in flight; 0: survival probability\\
\code{ctau} & cm & 0 & $c\tau$ of the hadron (\eqr{eq:survival})\\
\code{hadron\_arm} & code & 0 & also selects $z_\mathrm{aero}$ (\tabr{tab:zaero})\\
\code{pt\_b\_param} & GeV$^{-2}$ & 0 & read, not used (warning when set)\\
\code{sigc\_flag} & -- & 0 & read, not used (central-branch binning: 0 in $z$, 1 in $p_T^2$)\\
\code{sigc\_nbin} & -- & 10 & read, not used\\
\code{sigc\_kin\_min} & -- & 0 & read, not used\\
\code{sigc\_kin\_max} & -- & 1 & read, not used\\
\code{sigc\_kin\_ind} & -- & 0 & read, not used\\
\end{simctable}

\subsection{Outputs}
The default (upstream 59ccdf6) tree has 64 columns for SIDIS (76 with a $\pi^0$, 8 more with the polarised-target field);
\code{-{}-tree-layout legacy} gives the 56-column layout of SIMC 294f662 (68 with a $\pi^0$). Columns 1--40 are those of
the (e,e\pr p)-type layout (the arm blocks \code{hs}, \code{ss} and the kinematic columns, note 01, Outputs), with
\code{Em} of \eqr{eq:recon}, \code{Emi} and \code{Pmi} of \eqr{eq:pmiss} and \code{Wi}, \code{epsiloni} of the meson block.
The following columns are specific:

\begin{simctable}{footnotesize}{L{0.06\linewidth}L{0.14\linewidth}L{0.17\linewidth}L{0.63\linewidth}}
\caption{SIDIS tree columns 41--64 of the default layout (\code{FillUpstreamNtuple}; one column per row).}\lbl{tab:columns}\\
\toprule\rowcolor{simchead}\simcth{\#} & \simcth{name} & \simcth{unit} & \simcth{content}\\\midrule\endfirsthead
\toprule\rowcolor{simchead}\simcth{\#} & \simcth{name} & \simcth{unit} & \simcth{content}\\\midrule\endhead
\bottomrule\endlastfoot
41 & \code{missmass} & GeV & $M_X^\mathrm{rec}$ on a nucleon (\eqr{eq:recon})\\
42 & \code{mmnuc} & GeV & $M_{X,A}^\mathrm{rec}$ on the nucleus\\
43 & \code{phad} & GeV/$c$ & $P_h^\mathrm{rec}$\\
44 & \code{t} & GeV$^2$ & $t^\mathrm{rec}$ (positive convention)\\
45 & \code{ti} & GeV$^2$ & vertex $t$ (\eqr{eq:mesonkin})\\
46 & \code{z} & -- & $z^\mathrm{rec}$\\
47 & \code{zi} & -- & vertex $z$\\
48 & \code{pt2} & GeV$^2$ & $p_T^{2,\mathrm{rec}}$\\
49 & \code{pt2i} & GeV$^2$ & vertex $p_T^2$\\
50 & \code{thetapq} & rad & $\theta^\mathrm{rec}_{pq}$\\
51 & \code{phipq} & rad & $\varphi^\mathrm{rec}_{pq}$\\
52 & \code{thetapqi} & rad & vertex $\theta_{pq}$\\
53 & \code{phipqi} & rad & vertex $\varphi_{pq}$\\
54 & \code{pmpar} & GeV/$c$ & $P^\mathrm{rec}_{m,\parallel}$\\
55 & \code{pmper} & GeV/$c$ & $P^\mathrm{rec}_{m,\perp}$\\
56 & \code{pmoop} & GeV/$c$ & $P^\mathrm{rec}_{m,\mathrm{oop}}$\\
57 & \code{fry} & cm & minus the vertical raster coordinate\\
58 & \code{radphot} & GeV & $E_{\gamma1}+E_{\gamma2}+E_{\gamma3}$\\
59 & \code{pfermi} & GeV/$c$ & $p_F$ with the sign of $\hat p_F\cdot\hat q$ (0 for H and $A\ge3$; for deuterium the drawn $p_F$, also without \code{do\_fermi})\\
60 & \code{siglab} & $\mu$b/(MeV$^2$\,sr$^2$) & $\sigma$ of \eqr{eq:xsjac}, with the Coulomb factor\\
61 & \code{sighad} & GeV$^{-2}$ & the multiplicity factor of \eqr{eq:sigsemi} (0 below the cut of \eqr{eq:mmcut})\\
62 & \code{Weight} & $\mu$b/(MeV$^2$\,sr$^2$) & $w$ of \eqr{eq:weight}\\
63 & \code{decdist} & cm or -- & decay distance (\code{doing\_decay} = 1) or $P_\mathrm{surv}$ (\code{doing\_decay} = 0)\\
64 & \code{Mhadron} & MeV & mass of the detected particle (the decay product after a decay)\\
65--76 & photon columns & cm, MeV & $\pi^0$ only: calorimeter positions and four-momenta of the two photons\\
\end{simctable}
The legacy layout orders the columns differently and has some that the default layout does not: \code{xbj},
\code{xbji}, \code{thqi} (the vertex polar angle of $\vec q$), \code{jacobian} ($1000J_{zp_T}$), \code{sigcent} and
\code{centjac} (always 0) and \code{xfermi} (the $x$ of the model); its column list is in the reference manual
(\file{docs/manual/markdown/reference/OUTPUT\_FORMAT.md}, \S1e).

% =====================================================================================================================
\section{Validation in SIMC++}\lbl{sec:validation}
\begin{simctable}{footnotesize}{L{0.5\linewidth}L{0.21\linewidth}L{0.29\linewidth}}
\caption{CTest tests that cover SIDIS (one test per row).}\lbl{tab:tests}\\
\toprule\rowcolor{simchead}\simcth{test} & \simcth{labels} & \simcth{what it checks}\\\midrule\endfirsthead
\toprule\rowcolor{simchead}\simcth{test} & \simcth{labels} & \simcth{what it checks}\\\midrule\endhead
\bottomrule\endlastfoot
\code{l1\_sidis\_data\_files} & l1 & the data copies equal the Fortran files (md5)\\
\code{l1\_ctq5pdf} & l1 & \eqr{eq:cteq} against the Fortran table\\
\code{l1\_ctq5\_state} & l1 & the CTEQ5 table as read\\
\code{l1\_fdss} & l1 & \eqr{eq:fdss}\\
\code{l1\_f1f2in21} & l1 & \eqsr{eq:sfconv}{eq:resmod}\\
\code{l1\_slacemcfit\_rhad\_global} & l1 & \eqr{eq:rhad} and \eqr{eq:emc}\\
\code{l1\_peepix} & l1 & the complete model (compat mode)\\
\code{l1\_generate\_semi\_fixedrng} & l1 & the generation with fixed random numbers\\
\code{l1\_complete\_ev\_semi} & l1 & the vertex completion\\
\code{l1\_complete\_recon\_ev\_semi} & l1 & the reconstruction\\
\code{l1\_complete\_main\_semi} & l1 & the weight\\
\code{l1\_replay\_event\_capture\_semi\_h} & l1, replay & every routine of captured H events\\
\code{l1\_replay\_event\_capture\_semi\_d} & l1, replay & the same for deuterium\\
\code{l1\_montecarlo\_replay\_semi\_h} & l1, replay & the spectrometer chain of the H capture\\
\code{l1\_montecarlo\_replay\_semi\_d} & l1, replay & the same for deuterium\\
\code{l1\_simcrun\_replay\_semi\_h} & l1, replay & the complete run of the H capture, bitwise\\
\code{l1\_simcrun\_replay\_semi\_d} & l1, replay & the same for deuterium\\
\code{l1\_simcrun\_ranlux\_p14f2\_sidis\_h\_cap} & l1, replay & the $\pi^0$ capture with the calorimeter arm\\
\code{l1\_simcrun\_ranlux\_up\_sidis\_h\_k1\_pip\_rad\_A} & l1, replay, slow, upstream\_layout & full run, RANLUX and compat mode, against the upstream Fortran run\\
\code{cli\_sidis\_central\_log\_sidis\_h\_k1\_pip\_rad\_A} & cli, l1 & the central-event log lines\\
\code{cli\_sidis\_central\_log\_p14f2\_sidis\_h\_A} & cli, l1 & the same for the $\pi^0$ run\\
\code{sidis\_tables} & unit & table readers\\
\code{sidis\_sanity} & unit & limits, refusals and the default-mode corrections\\
\code{fortran\_diff\_D051\_sidis} & fortran\_diff, unit & the corrected model against the compat model\\
\code{fortran\_diff\_D051\_F3\_F9\_semi\_pfermi\_pmper} & fortran\_diff, unit & the \code{pfermi} column and the vertex $P_{m,\perp}$\\
\code{fortran\_diff\_D051\_F12\_F14\_semi\_hist\_titles} & fortran\_diff, unit & the \texttt{.hist} titles\\
\code{regression\_identity\_sidis\_h\_k1\_pip\_rad} & regression & compat-mode outputs against the stored baseline\\
\code{regression\_corrected\_sidis\_h\_k1\_pip\_rad} & regression & default-mode outputs against the stored baseline\\
\code{regression\_upstream\_u59\_sidis\_d\_k1\_pip\_fermi\_rad\_A} & regression\_upstream, slow & the deuterium run with Fermi motion against the upstream Fortran (which stops)\\
\code{regression\_shipped\_nps\_sidis\_pi0\_test} & regression\_shipped, slow & the shipped $\pi^0$ input\\
\end{simctable}
Recorded in the validation ledger (\code{project\_control/VALIDATION\_LEDGER.md}, entries of 2026-10-22 and 2026-10-24):
the model library is bitwise equal to the Fortran tables in the compat mode (43685 CTEQ5, 36432 fDSS, 360 F1F2IN21, 1696
nuclear-factor and 3338 complete-model rows), the default mode differing only where the corrections act; the statistical
comparison of the compat mode with the seven Fortran SIDIS runs (seed 4242, $10^5$ events; H, D, C-12, $\pi^\pm$, K$^+$, with
and without radiation) passes for all seven (8, 6, 8, 0, 5, 2 and 3 failed comparisons out of 253--270); the
default-mode differences are the documented corrections (the hadron $\delta$ cut of \code{pass\_cuts}, the survival
probability, the \code{pfermi} column).

% =====================================================================================================================
\section{Fortran origin and Fortran findings}\lbl{sec:origin}
\begin{simctable}{footnotesize}{L{0.42\linewidth}L{0.58\linewidth}}
\caption{Fortran routine $\to$ C++ function (one routine per row; the common routines are in \ctabr{tab:map}).}\lbl{tab:map}\\
\toprule\rowcolor{simchead}\simcth{Fortran (upstream 59ccdf6)} & \simcth{C++ (src/...)}\\\midrule\endfirsthead
\toprule\rowcolor{simchead}\simcth{Fortran (upstream 59ccdf6)} & \simcth{C++ (src/...)}\\\midrule\endhead
\bottomrule\endlastfoot
\file{dbase.f:125}-134, 192-207 (flags) & \csrc{core/RunSetup.cpp: DeriveReactionFlags}\\
\file{dbase.f:305}-317 (masses) & \csrc{setup/SetupReaction.cpp: BuildReactionAndKinematics}\\
\file{dbase.f:567}-591 (deut.dat) & \csrc{setup/SetupTheory.cpp: BuildTargetAndTheory, ReadMomentumDistribution}\\
\file{init.f:359}-513 (limits) & \csrc{setup/SetupLimits.cpp: BuildLimits}\\
\file{event.f} generate, l.276-373 & \csrc{reactions/sidis/SidisProduction.cpp: SidisProduction::generate}\\
\file{event.f} complete\_ev, l.723-786 & \csrc{reactions/pion/PionProduction.cpp: PionProduction::pionMainKinematics}\\
\file{event.f} complete\_ev, l.914-1047 & \csrc{reactions/sidis/SidisProduction.cpp: SidisProduction::completeEv}\\
\file{event.f} complete\_recon\_ev, l.1325-1341 & \csrc{reactions/EepReactionBase.cpp: EepReactionBase::completeRecon}\\
\file{event.f} complete\_main, l.1523-1564 & \csrc{reactions/sidis/SidisProduction.cpp: SidisProduction::completeMain}\\
\file{semi\_physics.f} peepiX & \csrc{physics/sidis/SidisCrossSection.cpp: peepiX}\\
\file{semi\_physics.f} (call of peepiX) & \csrc{physics/sidis/SidisModel.cpp: SidisModel::sigccMain}\\
\file{semi\_physics.f} slacemcfit & \csrc{physics/sidis/SidisCrossSection.cpp: slacemcfit}\\
\file{semi\_physics.f} Rhad\_global & \csrc{physics/sidis/SidisCrossSection.cpp: rhadGlobal}\\
\file{cteq5/Ctq5Pdf.f} SetCtq5, ReadTbl & \csrc{physics/sidis/Cteq5Pdf.cpp: Cteq5Pdf::load}\\
\file{cteq5/Ctq5Pdf.f} Ctq5Pdf & \csrc{physics/sidis/Cteq5Pdf.cpp: Cteq5Pdf::pdf}\\
\file{cteq5/Ctq5Pdf.f} PartonX & \csrc{physics/sidis/Cteq5Pdf.cpp: Cteq5Pdf::partonX}\\
\file{cteq5/Ctq5Pdf.f} POLINT & \csrc{physics/sidis/Cteq5Pdf.cpp: polint}\\
\file{fdss/fdss.f} fDSS (grid) & \csrc{physics/sidis/Fdss.cpp: FdssGrid::load}\\
\file{fdss/fdss.f} fDSS (evaluation) & \csrc{physics/sidis/Fdss.cpp: FdssGrid::eval}\\
\file{fdss/fdss.f} fFINT & \csrc{physics/sidis/Fdss.cpp: ffint}\\
\file{F1F2IN21\_v1.0.f} F1F2IN21 & \csrc{physics/sidis/F1F2In21.cpp: f1f2in21}\\
\file{F1F2IN21\_v1.0.f} SF & \csrc{physics/sidis/F1F2In21.cpp: sf}\\
\file{F1F2IN21\_v1.0.f} rescsp, RESCSN & \csrc{physics/sidis/F1F2In21.cpp: rescsp, rescsn}\\
\file{F1F2IN21\_v1.0.f} RESMODP, RESMODN & \csrc{physics/sidis/F1F2In21.cpp: resmod}\\
\file{F1F2IN21\_v1.0.f} SQESUB & not ported (no effect on SIDIS)\\
\file{semi\_dilution.f} semi\_dilution & \csrc{physics/sidis/SemiDilution.cpp: semiDilution}\\
\file{radc.f:297}-303 (tail-1 case) & \csrc{radiation/Radc.cpp: generateRad}\\
\file{simc.f:1298}-1313 (central log) & \csrc{driver/SimcRun.cpp}\\
\file{results\_write.f:178}-286 & \csrc{output/NtupleLayout.cpp: FillUpstreamNtuple}\\
\file{simc.f:394}-401 & \csrc{reactions/EepReactionBase.cpp: EepReactionBase::genVolume}\\
\end{simctable}

\begin{simcfindings}{footnotesize}{L{0.2\linewidth}L{0.5\linewidth}L{0.3\linewidth}}
\caption{Fortran findings that touch SIDIS (one finding per row; the common ones are in \ctabr{tab:findings}).}\lbl{tab:findings}\\
\toprule\rowcolor{simchead}\simcth{id} & \simcth{Fortran behaviour} & \simcth{default / compat}\\\midrule\endfirsthead
\toprule\rowcolor{simchead}\simcth{id} & \simcth{Fortran behaviour} & \simcth{default / compat}\\\midrule\endhead
\bottomrule\endlastfoot
\finding{D-051 F1} & all parton distributions 0: weight $0/0$ = NaN & 0 / NaN\\
\finding{D-051 F2} & negative Fermi $x$ stops the run in Ctq5Pdf & $\sigma=0$ / stop\\
\finding{D-051 F3} & \code{pfermi} 0/0 = NaN without Fermi momentum (fixed upstream) & 0 / 0\\
\finding{D-051 F4} & $\pi^0$: $D_{\bar d}$ not averaged, charged-pion mass in the model & averaged, $M_{\pi^0}$ / Fortran\\
\finding{D-051 F6 survival p} & survival probability with the central momentum & hadron momentum / Fortran\\
\finding{D-051 F6 zaero} & no aerogel distance for the SHMS; path without the square root & 0, square root / Fortran\\
\finding{D-051 F9} & vertex $P_{m,\perp}$ from the previous try's $P_m$ & current $P_m$ / Fortran\\
\finding{D-051 F12} & no $\pi^0$ title in the \texttt{.hist} & $\pi^0$ title / Fortran\\
\finding{D-051 F14} & $A\ge3$ kaon SIDIS stops when the \texttt{.hist} title is written & title written / stop\\
\finding{D-051 F7} & the ``2$\pi$/$\pi$K'' cut of \eqr{eq:mmcut} multiplies by $1-p_T^2$ & kept in both\\
\finding{D-051 F8} & threshold \eqr{eq:threshold} with $M_{\pi^0}$ for kaons, no Fermi motion; $z>1$ test unreachable & kept in both\\
\finding{D-051 F10} & $R_h\ne1$ for deuterium; EMC fit and $M_N=M_p$ choices; no Fermi momentum for $A\ge3$ & kept in both\\
\finding{D-051 F16} & hadron generated up to $\Sigma_\mathrm{max}-E'_\mathrm{min}$ with radiation & kept in both\\
\finding{D-051 F5 dead inputs} & \code{pt\_b\_param} read and not used & kept; warning\\
\finding{D-051 F5 doing\_cent} & central branch never called; \code{sigc\_*} inputs unused & ported, not called; warning\\
\finding{D-051 F5 semi\_dilution} & dilution routine never called & ported, not called\\
\finding{D-051 F13} & \code{CENTRAL.sigcc} from the reconstructed central event & kept\\
\finding{D-005 sigc\_nbin} & \code{sigc\_nbin} = 0 would divide by zero in the central branch & kept (branch not called)\\
\finding{P12.1 Q1} & early return of the model below the cut of \eqr{eq:mmcut} & kept\\
\finding{P12.1 Q2} & first call reads deuteron smearing tables (no effect) & not ported\\
\finding{P12.1 Q3} & a resonance-fit cut that is never true & kept\\
\finding{P12.1 Q4} & Ctq5Pdf overwrites the caller's $Q$ below $\Lambda$ & reproduced\\
\finding{P12.1 Q5} & fDSS arithmetic in quadruple precision & reproduced\\
\finding{P12.1 Q6} & pion survival probability written, not applied & kept\\
\finding{P12.1 Q7} & $x>1$ replaced by 1 with a message & kept\\
\finding{P12.1 Q8} & PartonX reads outside its arrays for $x\le$ the first node after the first call & reproduced\\
\finding{U-D1} & every SIDIS try with $z>1$ prints a warning (unreachable, \eqr{eq:zpt}) & followed\\
\finding{D-006 pass\_cuts} & hadron $\delta$ cut with the electron upper limit (\texttt{wtcontr} of the SIDIS runs $-28.9$ to $-32.2$\,\%) & hadron limit / Fortran\\
\finding{D-062 G8} & a tail-1 photon redraws the $\pi^0$ decay photons & kept\\
\finding{D-065 P1} & a $\pi^0$ transported through a magnetic arm & kept / stop with decay\\
\end{simcfindings}

\begin{simcfindingbox}
\textbf{Description of the early return.} The list of Fortran differences says that, below the cut of \eqr{eq:mmcut},
the columns \code{sighad} and \code{decdist} keep the previous event's values. Both are reset to 0 at the start of every
try (\fsrc{simc.f:184-186}, and in SIMC++ \csrc{driver/SimcRun.cpp}), so only the jacobian column keeps a stale value and
the survival probability is stale only as an internal variable (the weight is 0). Documentation only.
\end{simcfindingbox}

\begin{simcrefs}
\simcref{cteq5} H.~L.~Lai et al.\ (CTEQ), ``Global QCD analysis of parton structure of the nucleon: CTEQ5 parton
  distributions'', Eur.\ Phys.\ J.\ C \textbf{12}, 375 (2000), hep-ph/9903282 (cited at \fsrc{cteq5/Ctq5Pdf.f:5-8}).
\simcref{accardi} A.~Accardi, T.~Hobbs, W.~Melnitchouk, ``Hadron mass corrections in semi-inclusive deep inelastic
  scattering'', JHEP \textbf{11} (2009) 084, arXiv:0907.2395 (cited at \fsrc{semi\_physics.f:462}).
\simcref{dss} D.~de Florian, R.~Sassot, M.~Stratmann, Phys.\ Rev.\ D \textbf{75}, 114010 (2007) and D \textbf{76},
  074033 (2007) (cited at \fsrc{fdss/fdss.f:4-5}).
\simcref{christy} M.~E.~Christy et al., F1F2IN21 fit, version 1.0 prerelease of April 4, 2022 (named at
  \fsrc{F1F2IN21\_v1.0.f:3-20}; unpublished at that date according to the code).
\simcref{gomez} J.~Gomez et al., ``Measurement of the A dependence of deep-inelastic electron scattering'', Phys.\ Rev.\ D
  \textbf{49}, 4348 (1994): the usual source of the SLAC EMC fit; not cited in the code, not compared.
\simcref{binnewies} J.~Binnewies, B.~A.~Kniehl, G.~Kramer, Phys.\ Rev.\ D \textbf{52}, 4947 (1995) (the parametrisation of
  the unused dilution routine, cited at \fsrc{semi\_physics.f:14-15}).
\end{simcrefs}
\end{document}
